\section{Le théorème de dualité globale}\etiquette[3]{XVIII.sec3}\pageoriginale 

\subsection{Le foncteur $Rf^{!}$}\etiquette[3.1]{XVIII.sec3.1}

Nous aurons besoin d'une variante des résolutions flasques canoniques.

\begin{enonce*}{Lemme 3.1.1}\etiquette[3.1.1]{XVIII.lem3.1.1}
Soit $X$ un schéma cohérent (\ie quasi-compact quasi-séparé). Le
foncteur $\varinjlim$ induit une équivalence entre la catégorie des
Ind-objets de la catégorie des faisceaux d'ensembles (\resp de
groupes, \resp de groupes abéliens) constructibles et la catégorie des
faisceaux d'ensembles (\resp de groupes Ind-finis, \resp abéliens de torsion).
\end{enonce*}

Résulte de \hyperrefext[IX][IX.cor2.7.2]{2.7.2} et \hyperrefext[IX][IX.cor2.7.3]{2.7.3}.

\setcounter{subsubsection}{1}
\subsubsection{}\etiquette[3.1.2]{XVIII.subsubsec3.1.2}
Soit $\Ff$ un faisceau abélien de torsion sur un schéma $X$,
limite inductive filtrante de faisceaux constructibles
$(\Ff_{i})_{i\in I}$, et $X_{0}$ un ensemble conservatif de
points de $X$. On définit la \emph{résolution flasque canonique
modifiée} de $\mathcal{F}$ par la formule
$$
C^{\ast}_{\ell}(\Ff)=\varinjlim_{i}C^{\ast}(\Ff_{i}){\quoi},
$$ où les $C^{\ast}(\Ff_{i})$ sont les résolutions flasques canoniques
de \hyperrefext[XVII][XVII.def4.2.2]{4.2.2}. D'après \hyperref[XVIII.lem3.1.1]{3.1.1}, pour $X$
cohérent, cette définition ne dépend pas du choix des $\Ff_{i}$. Elle
commute de plus à la localisation; ce fait permet de définir
$C^{\ast}_{\ell}(\Ff)$ par globalisation pour $\Ff$ faisceau de
torsion sur un schéma $X$ quelconque, non nécessairement cohérent.

\emph{Le complexe $C^{\ast}_{\ell}(\Ff)$ est une résolution de
$\Ff$, fonctorielle en $\Ff$, dont la formation
commute à la localisation et aux limites inductives filtrantes; les
foncteurs $C^{n}_{\ell}$ sont exacts.}

Il en résulte que pour un faisceau d'anneaux $\mathcal{A}$ sur $X$,
les foncteurs $C^{n}_{\ell}$ transforment $\mathcal{A}$-Modules en
$\mathcal{A}$-Modules; ils transforment de même faisceaux de torsion
en faisceaux de torsion. Il est bien connu que 

\begin{enonce*}{Lemme 3.1.3}\etiquette[3.1.3]{XVIII.lem3.1.3}
Soit\pageoriginale $\mathcal{A}$ une $\sheaf{\mathcal{U}}$-catégorie
abélienne ayant un petit ensemble générateur et telle que les limites
inductives indexées par un petit ensemble ordonné filtrant soient
représentables dans $\mathcal{A}$, et soient exactes. Pour qu'un
foncteur $F$ de $\mathcal{A}^{\circ}$ dans $(\Ens)$ soit représentable, il
faut et il suffit qu'il transforme petites limites inductives
(quelconques) en limites projectives.
\end{enonce*}

\Cf p.~ex. \SGA 4 I 8.12.5, où on ne suppose pas $\mathcal{A}$
abélienne ni les $\varinjlim$ filtrantes exactes.

On appliquera ce lemme à la catégorie $\Mod(X,\mathcal{A})$ des
faisceaux de modules sur un site annelé $(X,\mathcal{A})$. Dans ce cas
particulier, si le foncteur $F$ est représenté par le faisceau
$ℱ$, on a
\begin{equation*}
ℱ(U)=F(\mathcal{A}_{U}){\quoi}.\tag{3.1.3.1}
\end{equation*}\etiquette[3.1.3.1]{XVIII.eq3.1.3.1}

Cette description de $ℱ$ permet une vérification directe de
\hyperref[XVIII.lem3.1.3]{3.1.3} dans le cas considéré.

Rappelons que si $f:X\rightarrow S$ est un morphisme compactifiable
(\hyperrefext[XVII][XVII.subsec3.2]{3.2}.), alors $S$ est quasi-compact quasi-séparé (sic), et la
dimension des fibres de $f$ est bornée.

\begin{enonce*}{Théorème 3.1.4}\etiquette[3.1.4]{XVIII.thm3.1.4}
Soient $f:X\rightarrow S$ un morphisme compactifiable et $\mathcal{A}$
un faisceau d'anneaux de torsion sur $S$. Alors, le foncteur
\begin{equation*}
Rf_{!}:D(X,f^{\ast}\mathcal{A})\longrightarrow D(S,\mathcal{A})\tag{3.1.4.1}
\end{equation*}\etiquette[3.1.4.1]{XVIII.eq3.1.4.1}%
admet un adjoint à droite partiel
\begin{equation*}
Rf^{!}:D^{+}(S,\mathcal{A})\longrightarrow D^{+}(X,f^{\ast}\mathcal{A}):\tag{3.1.4.2}
\end{equation*}\etiquette[3.1.4.2]{XVIII.eq3.1.4.2}%
pour $K\in \Ob D(X,f^{\ast}\mathcal{A})$ et $L\in \Ob
D^{+}(S,\mathcal{A})$, on a un isomorphisme fonctoriel 
\begin{equation*}
\Hom(Rf_{!}K,L)\simeq \Hom(K,Rf^{!}L){\quoi}.\tag{3.1.4.3}
\end{equation*}\etiquette[3.1.4.3]{XVIII.eq3.1.4.3}%
Muni du morphisme de translation rendant commutatif le diagramme 
\begin{equation*}
\vcenter{
\xymatrix{
  \Hom(Rf_!(K[1]),L)
  \ar@{<->}[d]_{\text{\rotatebox{90}{$\backsim$}}}\ar@{<->}[r]^-{\sim}&
  \Hom(K[1],Rf^{!}L)\ar@{-}[r]^-\sim&
\Hom(K,(Rf^{!}L)[-1])\ar@{<->}[d]^{\text{\rotatebox{90}{$\backsim$}}}\\
\Hom((Rf_!K)[1],L)\ar@{-}[r]^-\sim & \Hom(Rf_!K,L[-1]
\ar@{<->}[r]^-{\sim}_-{\txt{\rm adj.}} & \Hom(K,Rf^{!}(L[-1])),
}}\tag{3.1.4.4}
\end{equation*}\etiquette[3.1.4.4]{XVIII.eq3.1.4.4}%
ce\pageoriginale foncteur $Rf^{!}$ est un foncteur triangulé.
\end{enonce*}

N.B. La notation $Rf^{!}$ est abusive en ce que $Rf^{!}$ n'est en
général pas le dérivé d'un foncteur $f^{!}$.

\noindent
\emph{Démonstration} Soit $d$ un entier tel que toute fibre de $f$
soit de dimension $<d$, et choisissons une compactification
$$
\xymatrix@R=1.25cm@C=1.5cm{
X \ar@{^{(}->}[r]^{j} \ar[d]_{f}& \bar{X}\ar[ld]^{\bar{f}}\\
S.
}
$$

Si $K$ est un complexe de $f^{\ast}\mathcal{A}$-Modules sur $X$, on a
$$
Rf_{!}K\simeq R\bar{f}_{\ast}(j_{!}K){\quoi}.
$$

Pour tout faisceau $F$ sur $X$, les composants $F^{i}$ du complexe
$\tau_{\leq 2d}C^\ast_{\ell}j_{!}F$ vérifient
\begin{equation*}
R^{k}\bar{f}_{\ast}F^{i}=0\quad \text{\rm pour}\quad k>0{\quoi}.\tag{3.1.4.5}
\end{equation*}\etiquette[3.1.4.5]{XVIII.eq3.1.4.5}%
En effet, pour $i\neq 2d$, $F^{i}$ est limite inductive de faisceaux
flasques, et pour $i=2d$ et $k>0$, on a (\hyperrefext[XVII][XVII.cor5.2.8.1]{5.2.8.1})
$$
R^{k}\bar{f}_{\ast}F^{2d}=R^{k+2d}f_{!}F=0{\quoi}.
$$

Le complexe simple associé au double complexe résolution flasque
canonique modifiée tronquée de $K$ est une résolution de $K$. En vertu
de \eqhyperref[XVIII.eq3.1.4.5]{3.1.4.5}, on a donc (notation de \hyperrefext[XVII][XVII.subsubsec1.1.15]{1.1.15})
\begin{equation*}
R\bar{f}_{\ast}(j_{!}K)\xleftarrow{\sim} \bar{f}_{\ast}\tau''_{\leq
2d} C^{\ast}_{\ell}j_{!}K{\quoi}.\tag{3.1.4.6}
\end{equation*}\etiquette[3.1.4.6]{XVIII.eq3.1.4.6}

Désignons par $f^{\bigdot}_{!}$ le foncteur des faisceaux de modules
sur $X$ dans les\pageoriginale complexes de faisceaux sur $S$ défini
par
\begin{equation*}
f^{\bigdot}_{!}(F)=\bar{f}_{\ast}\tau_{\leq
2d}C^{\ast}_{\ell}j_{!}F{\quoi}.\tag{3.1.4.7}
\end{equation*}\etiquette[3.1.4.7]{XVIII.eq3.1.4.7}

\begin{enonce*}{Lemme 3.1.4.8}\etiquette[3.1.4.8]{XVIII.lem3.1.4.8}
\begin{enumerate}
\renewcommand{\theenumi}{\roman{enumi}}
\renewcommand{\labelenumi}{(\theenumi)}
\item Les foncteurs $f^{i}_{!}$ sont exacts et commutent aux limites
inductives filtrantes; on a $f^{i}_{!}=0$ pour $i\not\in[0,2d]$.

\item Le foncteur $f^{\bigdot}_{!}$, étant borné à composantes
exactes, se dérive \emph{(\hyperrefext[XVII][XVII.prop1.2.10]{1.2.10})} trivialement en
$Rf^{\bigdot}_{!}$ ou simplement
$f^{\bigdot}_{!}:D(X,f^{\ast}\mathcal{A})\rightarrow
D(S,\mathcal{A})$. La flèche \eqhyperref[XVIII.eq3.1.4.6]{3.1.4.6} est un isomorphisme de
foncteurs
$$
Rf^{\bigdot}_{!}\xrightarrow{\sim} Rf_{!}{\quoi}.
$$
\end{enumerate}
\end{enonce*}

Les foncteurs $C^{i}_{\ell}$ sont exacts et commutent aux limites
inductives filtrantes (\hyperref[XVIII.subsubsec3.1.2]{3.1.2}). Ils transforment
tout faisceau en un faisceau acyclique. Les foncteurs $\bar{f}_{\ast}$
et $j_{!}$ commutent aux limites inductives filtrantes. On en déduit que
\begin{enumerate}
\renewcommand{\theenumi}{\alph{enumi}}
\renewcommand{\labelenumi}{\theenumi)}
\item les foncteurs $\bar{f}_{\ast}C^{i}_{\ell}j_{!}$, donc aussi les
foncteurs $f^{i}_{!}$, commutent aux limites inductives filtrantes;

\item pour $i<2d$, les foncteurs
$f^{i}_{!}=\bar{f}_{\ast}C^{i}_{\ell}j_{!}$ sont exacts. Pour $i=2d$,
on vérifie que $f^{i}_{!}$ est exact par \eqhyperref[XVIII.eq3.1.4.5]{3.1.4.5}.
\end{enumerate}

Les assertions restantes de \hyperref[XVIII.lem3.1.4.8]{3.1.4.8} sont triviales.

\begin{enonce*}[remark]{Remarque 3.1.4.9}\etiquette[3.1.4.9]{XVIII.rem3.1.4.9}
La définition de $f^{\bigdot}_{!}$ est indépendante du faisceau
d'anneau $\mathcal{A}$, et garde un sens pour tout faisceau
abélien. L'assertion \hyperref[XVIII.lem3.1.4.8]{3.1.4.8} (i) reste valable dans les
catégories de faisceaux abéliens de torsion sur $X$ et $S$.
\end{enonce*}

D'après \hyperref[XVIII.lem3.1.4.8]{3.1.4.8} et \hyperref[XVIII.lem3.1.3]{3.1.3}, le foncteur
$$
f^{i}_{!}:\Mod(X,f^{\ast}\mathcal{A})\longrightarrow
\Mod(S,\mathcal{A})
$$
a un adjoint à droite $f^{!}_{i}$; puisque $f^{i}_{!}$ est exact,
$f^{!}_{i}$ transforme injectifs en injectifs.

Avec la convention de signe (\hyperrefext[XVII][XVII.subsubsec1.1.12]{1.1.12}), les foncteurs $f^{!}_{i}$
forment un complexe de foncteurs exacts à gauche, nul en degré
cohomologique $n\not\in [-2d,0]$:
\begin{equation*}
f^{!}_{\bigdot}:\Mod(S,\mathcal{A})\longrightarrow
C^{b}(X,f^{\ast}\mathcal{A}){\quoi}.\tag{3.1.4.10}
\end{equation*}\etiquette[3.1.4.10]{XVIII.eq3.1.4.10}

Le\pageoriginale foncteur $f^{!}_{\bigdot}$ se dérive
(\hyperrefext[XVII][XVII.prop1.2.10]{1.2.10}) en un foncteur triangulé
\begin{equation*}
Rf^{!}:D^{+}(S,\mathcal{A})\longrightarrow
D^{+}(X,f^{\ast}\mathcal{A}){\quoi},\tag{3.1.4.11}
\end{equation*}\etiquette[3.1.4.11]{XVIII.eq3.1.4.11}%
adjoint à droite au foncteur $Rf_{!}$
$(=Rf^{\bigdot}_{!}=Lf^{\bigdot}_{!})$ (« formule triviale de
dualité »). Plus précisément, si $I\in \Ob C^{+}(S,\mathcal{A})$
est un complexe de faisceaux injectifs de $\mathcal{A}$-modules, les
$\mathcal{A}$-modules $f^{!}_{i} I^{n}$ sont injectifs, et on dispose
d'un isomorphisme de triples complexes (\hyperrefext[XVII][XVII.subsubsec1.1.12]{1.1.12})
\begin{equation*}
\Hom^{\bigdot}(f^{\bigdot}_{!}K, I)\xrightarrow{\sim}
\Hom^{\bigdot}(K,f^{!}_{\bigdot}I){\quoi}.\tag{3.1.4.12}
\end{equation*}\etiquette[3.1.4.12]{XVIII.eq3.1.4.12}%
Passant au $H^{0}$ du complexe simple associé (calculé par produits),
on obtient l'énoncé d'adjonction. On laisse au lecteur le soin de
vérifier la compatibilité \eqhyperref[XVIII.eq3.1.4.4]{3.1.4.4}.

\begin{enonce*}[remark]{Remarque 3.1.5}\etiquette[3.1.5]{XVIII.rem3.1.5}
Supposons que le foncteur $Rf^{!}$ soit d'amplitude cohomologique
finie. D'après \hyperrefext[XVII][XVII.prop1.2.10]{1.2.10} le foncteur $f^{!}_{\bigdot}$ admet alors
un dérivé
\begin{equation*}
Rf^{!}:D(S,\mathcal{A})\longrightarrow D(X,\mathcal{A}){\quoi}.\tag{3.1.5.1}
\end{equation*}\etiquette[3.1.5.1]{XVIII.eq3.1.5.1}%
Ce dérivé est encore adjoint à droite au foncteur $Rf_{!}$
\eqhyperref[XVIII.eq3.1.4.1]{3.1.4.1}. En effet, on a
\begin{align*}
& \Hom_{D(S)}(Rf_{!}K,L)\simeq
\varinjlim_{L\xrightarrow{\sim}L'}\Hom_{K(S)}(f^{\bigdot}_{!}K,L')\simeq
\varinjlim_{\substack{K'\xrightarrow{\sim}K\\ L\xrightarrow{\sim}
L'}}\Hom_{K(S)} (f^{\bigdot}_{!}K', L')\\
& \Hom_{D(X)}(K,Rf^{!}L)\simeq
\varinjlim_{K'\xrightarrow{\sim}K}\Hom_{K(X)}(K',Rf^{!}L)\simeq
\varinjlim_{\substack{K'\xrightarrow{\sim}K\\
L\xrightarrow{\sim}L'}}\Hom_{K(X)}(K', f^{!}_{\bigdot}L')
\end{align*}
et, par adjonction,
$$
\Hom_{K(S)}(f^{\bigdot}_{!}K',L')=H^{0}\Hom(f^{\bigdot}_{!}K',L')\simeq
H^{0}\Hom(K',f^{!}_{\bigdot}L')=\Hom_{K(X)}(K',L'){\quoi}.
$$
\end{enonce*}

\begin{enonce*}{Définition 3.1.6}\etiquette[3.1.6]{XVIII.def3.1.6}
Le foncteur triangulé
$$
Rf^{!}:D^{+}(S,\mathcal{A})\longrightarrow
D^{+}(X,f^{\ast}\mathcal{A})
$$
(ou,\pageoriginale sous les hypothèses de \hyperref[XVIII.rem3.1.5]{3.1.5},
$$
Rf^{!}:D(S,\mathcal{A})\longrightarrow
D(X,f^{\ast}\mathcal{A})){\quoi},
$$
adjoint à droite au foncteur $Rf_{!}$ \eqhyperref[XVIII.eq3.1.4.1]{3.1.4.1} s'appelle le
\emph{foncteur image inverse extraordinaire.}
\end{enonce*}

\begin{enonce*}{Proposition 3.1.7}\etiquette[3.1.7]{XVIII.prop3.1.7}
Supposons le morphisme compactifiable $f:X\rightarrow S$ de dimension
relative $\leq d$.
\begin{enumerate}
\renewcommand{\theenumi}{\roman{enumi}}
\renewcommand{\labelenumi}{(\theenumi)}
\item Si $L\in \Ob D^{+}(S,\mathcal{A})$ vérifie $H^{i}(L)=0$ pour
$i\leq k$, alors $H^{i}Rf^{!}L=0$ pour $i\leq k-2d$.

\item Pour $L\in \Ob D^{b}(S,\mathcal{A})$, on a
$$
\dim \inj (Rf^{!}L)\leq \dim \inj (L){\quoi}.
$$
\end{enumerate}
\end{enonce*}

Le foncteur $Rf^{!}$ est le dérivé droit de
$f^{!}_{\bigdot}:\Mod(S,\mathcal{A})\rightarrow
C^{b}\Mod(X,f^{\ast}\mathcal{A})$. L'assertion (i) résulte de ce que
$(f^{!}_{\bigdot})^{i}=0$ pour $i<-2d$, et l'assertion (ii) résulte de
ce que $(f^{!}_{\bigdot})^{i}=0$ pour $i>0$, et que les foncteurs
$f^{!}_{i}$ transforment injectif en injectif.

\begin{enonce*}{Proposition 3.1.8}\etiquette[3.1.8]{XVIII.prop3.1.8}
Soit $f:X\rightarrow S$ un morphisme compactifiable quasi-fini.
\begin{enumerate}
\renewcommand{\theenumi}{\roman{enumi}}
\renewcommand{\labelenumi}{(\theenumi)}
\item Le foncteur $f_{!}:\Mod(X,f^{\ast}\mathcal{A})\rightarrow
\Mod(S,\mathcal{A})$ admet un adjoint à droite
$f^{!}:\Mod(S,\mathcal{A})\rightarrow \Mod(X,f^{\ast}\mathcal{A})$, et
$Rf^{!}$ est le foncteur dérivé du foncteur $f^{!}$.

\item Si $f$ est une immersion fermée, alors $f^{!}$ est le foncteur
« sections à support dans $X$ ».

\item Si $f$ est étale, alors $f^{!}$ s'identifie au foncteur
$f^{\ast}$, avec 
$$
\Tr_{f}:f_{!}f^{\ast}F\longrightarrow F
$$
comme morphisme d'adjonction.
\end{enumerate}
\end{enonce*}

L'assertion (i) résulte aussitôt du cas particulier $d=0$ de la
démonstration de \hyperref[XVIII.thm3.1.4]{3.1.4}: on a $f^{\bigdot}_{!}=f_{!}$ et
$f^{!}_{\bigdot}=f^{!}$. De façon équivalente, l'existence d'un
foncteur adjoint $f^{!}$ résulte de \hyperref[XVIII.lem3.1.3]{3.1.3} et de ce que $f_{!}$
est\pageoriginale exact et commute aux limites inductives filtrantes;
l'adjonction $(f_{!},Rf^{!})$ est une « formule triviale de
dualité ».

L'assertion (ii) est claire, et (iii) est \hyperrefext[XVII][XVII.prop6.2.11]{6.2.11}.

\setcounter{subsubsection}{8}
\subsubsection{}\etiquette[3.1.9]{XVIII.subsubsec3.1.9}
Soient $f:X\rightarrow S$ un morphisme compactifiable, $\mathcal{A}$
un faisceau d'anneaux de torsion sur $S$, et annelons $X$ par l'image
réciproque de $\mathcal{A}$. Soient $K\in \Ob
D^{-}(X,f^{\ast}\mathcal{A})$, $L\in \Ob D^{+}(X,f^{\ast}\mathcal{A})$
et représentons $L$ par un complexe borné inférieurement de faisceaux
injectifs, de sorte que
$$
R\Hhom(K,L)\simeq \Hhom(K,L){\quoi}.
$$
Le complexe $\Hhom(K,L)$ est à composantes flasques, et donc
\begin{equation*}
Rf_{\ast}R\Hhom(K,L)\simeq f_{\ast}\Hhom(K,L){\quoi}.\tag{3.1.9.1}
\end{equation*}\etiquette[3.1.9.1]{XVIII.eq3.1.9.1}%
Choisissons une compactification de $f$ et un entier $d$ comme dans la
démonstration de \hyperref[XVIII.thm3.1.4]{3.1.4}, d'où un foncteur
$f^{\bigdot}_{!}$ \eqhyperref[XVIII.eq3.1.4.7]{3.1.4.7}. 
\begin{equation*}
\vcenter{
\xymatrix@C=1.5cm@R=.5cm{
X_V \ar[rr]^{k_X} \ar@{^{(}->}[rd]_{j_V} \ar[ddd]^{f_V}&      & X\ar[ddd]^{f} \ar@{^{(}->}[rd]^{j}         \\
    & \bar{X}_V \ar[ldd]^{\bar{f}_V} \ar[rr]_(.3){k_{\bar{X}}}   &   & \bar{X} \ar[ldd]_{\bar{f}}\\
    &         &        &\\
V  \ar[rr]^{k} &      & S.
}}\tag{3.1.9.2}
\end{equation*}\etiquette[3.1.9.2]{XVIII.eq3.1.9.2}


Pour tout $V\in \Ob S_{\text{\rm et}}$, on a un isomorphisme\nde{Ce
morphisme est donné par la construction de \hyperrefext[XVII][XVII.rem4.2.6]{4.2.6}, mais il n'y
a aucune raison que ce soit un isomorphisme si on n'a pas préalablement
fait des choix cohérents de points sur les topos considérés. Pour que ce
soit le cas, il faut que le
carré qui intervient dans \hyperrefext[XVII][XVII.rem4.2.6]{4.2.6} soit en quelque sorte
cartésien. Plus précisément, il convient ici de fixer une famille de points
$(p)_{p\in P}$
du topos $S_{\text{\rm et}}$ et pour tout $V\in \Ob S_{\text{\rm et}}$, on
considère la résolution flasque canonique modifiée utilisant la famille de
points associés par la correspondance de \hyperrefext[IV][IV.subsec6.7]{6.7}
aux couples $(p,\xi)$ où $p\in P$ et $\xi\in V_p$.}
\begin{equation*}
k^{\ast}f^{\bigdot}_{!}\xrightarrow{\sim}
f^{\bigdot}_{V!}k^{\ast}_{X}{\quoi},\tag{3.1.9.3}
\end{equation*}\etiquette[3.1.9.3]{XVIII.eq3.1.9.3}%
de sorte que
\begin{equation*}
\Hhom^{\bigdot}(f^{\bigdot}_{!}K,f^{\bigdot}_{!}L)(V)\simeq
\Hhom^{\bigdot}(f^{\bigdot}_{V!}(k^{\ast}_{X}K),f^{\bigdot}_{V!}(k^{\ast}_{X}L)){\quoi}.
\end{equation*}%
La fonctorialité de $f^{\bigdot}_{V!}$ nous fournit
$$
\Hhom^{\bigdot}(K,L)(X_{V})=\Hom^{\bigdot}(k^{\ast}_{X}K,k^{\ast}_{X}L)\longrightarrow
\Hom^{\bigdot}(f^{\bigdot}_{V!}(k^{\ast}_{X}K),f^{\bigdot}_{V!}(k^{\ast}_{X}L)){\quoi},
$$
d'où une flèche, au niveau des complexes,
\begin{equation*} 
f_{\ast}\Hhom^{\bigdot}(K,L)\longrightarrow
\Hhom^{\bigdot}(f^{\bigdot}_{!}K,f^{\bigdot}_{!}L){\quoi}.\tag{3.1.9.4} 
\end{equation*}\etiquette[3.1.9.4]{XVIII.eq3.1.9.4}

D'après\pageoriginale \eqhyperref[XVIII.eq3.1.9.1]{3.1.9.1}, cette flèche se dérive en
\begin{equation*}
Rf_{\ast}R\Hhom(K,L)\longrightarrow
R\Hhom(Rf_{!}K,Rf_{!}L){\quoi}.\tag{3.1.9.5}
\end{equation*}\etiquette[3.1.9.5]{XVIII.eq3.1.9.5}

Soient maintenant $K\in \Ob D^{-}(X,\mathcal{A})$ et $L\in \Ob
D^{+}(S,\mathcal{A})$. La flèche d'adjonction de $Rf_{!}Rf^{!}L$ dans
$L$ et \eqhyperref[XVIII.eq3.1.9.4]{3.1.9.4} fournissent une flèche composée
$$
Rf_{\ast}R\Hhom(K,Rf^{!}L)\longrightarrow R\Hhom(Rf_{!}K,
Rf_{!}Rf^{!}L)\longrightarrow R\Hhom(Rf_{!}K,L){\quoi},
$$
forme locale sur $S$ de l'isomorphisme d'adjonction:
\begin{equation*}
Rf_{\ast}R\Hhom(K,Rf^{!}L)\longrightarrow
R\Hhom(Rf_{!}K,L){\quoi}.\tag{3.1.9.6}
\end{equation*}\etiquette[3.1.9.6]{XVIII.eq3.1.9.6}

\begin{enonce*}{Proposition 3.1.10}\etiquette[3.1.10]{XVIII.prop3.1.10}
La flèche \hyperref[XVIII.eq3.1.9.6]{3.1.9.6} est un isomorphisme.
\end{enonce*}

Soit $V\in \Ob S_{\text{\rm et}}$ et reprenons les notations de
\eqhyperref[XVIII.eq3.1.9.2]{3.1.9.2}. L'isomorphisme (« de transitivité »)
$$
k_{!}Rf_{V!}\simeq Rf_{!}k_{X!}{\quoi},
$$
dont l'existence est la racine de la théorie du foncteur $Rf_{!}$, se
transpose en un \emph{isomorphisme de localisation} (\hyperref[XVIII.prop3.1.8]{3.1.8}. (iii))
\begin{equation*}
k^{\ast}_{X}Rf^{!}\simeq Rf^{!}_{V}k^{\ast}{\quoi},\tag{3.1.10.1}
\end{equation*}\etiquette[3.1.10.1]{XVIII.eq3.1.10.1}%
représenté au niveau des complexes par un morphisme
$$
k^{\ast}_{X}f^{!}_{\bigdot}\longrightarrow f^{!}_{V\bigdot}k^{\ast}{\quoi}.
$$

Pour vérifier que \eqhyperref[XVIII.eq3.1.9.6]{3.1.9.6} est un isomorphisme, il suffit de vérifier
que pour tout $V\in \Ob S_{\text{\rm et}}$, \eqhyperref[XVIII.eq3.1.9.6]{3.1.9.6} induit un
isomorphisme 
\begin{align*}
\Hom &(k^{\ast}_{X}K,k^{\ast}_{X}Rf^{!}L)\simeq
 H^{0}(V,Rf_{\ast}R\Hhom (K,Rf^{!}L))\tag{3.1.10.2}\\
& \longrightarrow H^{0}(V,R\Hom(Rf_{!}K,L))\simeq
 \Hhom(k^{\ast}Rf_{!}K,k^{\ast}L)
\end{align*}\etiquette[3.1.10.2]{XVIII.eq3.1.10.2}

Pour $V=X$, ceci n'est autre que \eqhyperref[XVIII.eq3.1.4.3]{3.1.4.3}. Le cas général en résulte
via la compatibilité suivante, \emph{que le rédacteur n'a pas vérifiée}:

\begin{enonce*}{Lemme 3.1.10.3}\etiquette[3.1.10.3]{XVIII.lem3.1.10.3}
Le\pageoriginale diagramme
$$
\xymatrix@C=2.5cm@R=1.5cm{
\Hom(k^{\ast}_XK, k^{\ast}_XRf^{!}L)
\ar[r]^{\scriptstyle{(3.1.10.2)}} \ar[d]^{\scriptstyle{(3.1.10.1)}}_{\text{\rotatebox{90}{$\backsim$}}}&
\Hom(k^{\ast}Rf_!K,k^{\ast}L) \ar[d]^{\txt{\rm (localisation)}}_{\text{\rotatebox{90}{$\backsim$}}}\\
\Hom(k^{\ast}_X,K, Rf^{!}_Vk^{\ast}_XL)
\ar@{<->}[r]^{\scriptstyle{\rm{(3.1.4.3\; sur\; V)}}}_{\sim}& \Hom(Rf_!k^{\ast}K,k^{\ast}L)
}
$$
est commutatif.
\end{enonce*}

\setcounter{subsubsection}{10}
\subsubsection{}\etiquette[3.1.11]{XVIII.subsubsec3.1.11}
Soient un carré cartésien
\begin{equation*}
\vcenter{
\xymatrix@C=2cm@R=1.5cm{
X' \ar[r]^{g'}\ar[d]^{f'}& X\ar[d]^{f}\\
S' \ar[r]^{g} & S,
}}\tag{3.1.11.1}
\end{equation*}\etiquette[3.1.11.1]{XVIII.eq3.1.11.1}%
avec $f$ compactifiable, $\mathcal{A}$ un faisceau d'anneaux de
torsion sur $S$, $\mathcal{B}$ un faisceau d'anneaux de torsion sur
$S'$ et $M$ un complexe de
$(f^{\ast}\mathcal{A},\mathcal{B})$-bimodules, borné supérieurement. On dispose alors, au niveau des
complexes, d'un morphisme de foncteurs en $K$ $(K\in \Ob
C(X,f^{\ast}\mathcal{A}))$
\begin{equation*}
M\otimes_{\mathcal{A}}g^{\ast}f^{\bigdot}_{!}K\longrightarrow
f'_{!}{}^{\bigdot}
({f'}^{\ast}M\otimes_{\mathcal{A}}{g'}^{\ast}K){\quoi},\tag{3.1.11.2}
\end{equation*}\etiquette[3.1.11.2]{XVIII.eq3.1.11.2}%
qui se transpose en un morphisme de foncteurs en $L$ $(L\in \Ob
C(S',\mathcal{B}))$
\begin{equation*}
g'_{\ast}\Hhom_{\mathcal{B}}({f'}^{\ast}M,f'_{\bigdot}{}^{!}L)\longrightarrow
f^{!}_{\bigdot}g_\ast\Hhom_{\mathcal{B}}(M,L){\quoi}.\tag{3.1.11.3}
\end{equation*}\etiquette[3.1.11.3]{XVIII.eq3.1.11.3}

Ce morphisme se dérive en morphisme de foncteurs en $L$, de $D^{+}$
$(S',\mathcal{B})$ dans $D^{+}(X,f^{\ast}\mathcal{A})$:
\begin{equation*}
Rg'_{\ast}R\Hhom_{\mathcal{B}}({f'}^{\ast}M,Rf'{}^{!}L)\longrightarrow
Rf^{!}Rg_{\ast}R\Hom_{\mathcal{B}}(M,L){\quoi},\tag{3.1.11.4}
\end{equation*}\etiquette[3.1.11.4]{XVIII.eq3.1.11.4}%
qui rend commutatif le diagramme suivant, où $K\in \Ob
D^{-}(X,f^{\ast}\mathcal{A})$ et $L\in \Ob D^{+}(S',\mathcal{B})$:
$$
\xymatrix{
\scriptscriptstyle{\Hom(K,Rg'_{\ast} R\Hhom_{\Bb}(f'^{\ast}M,Rf'^{!}L))}
\ar@{<->}[d] \ar[r]^{\text{\ding{172}}} &
\scriptscriptstyle{\Hom(K,Rf^{!}Rg_{\ast}R\Hhom_{\Bb}(M,L))}
\ar@{<->}[r]_{{\scriptscriptstyle{\rm adj}}}& \scriptscriptstyle{\Hom(Rf_!K, Rg_{\ast}
R\Hhom_{\Bb}(M,L))}\ar@{<->}[d]\\
\scriptscriptstyle{\Hom(g'^{\ast}K,R\Hhom_{\Bb}(f'^{\ast} M, Rf'^{!}L)}\ar@{<->}[d] 
& &  \scriptscriptstyle{\Hom(g^{\ast}Rf_! K, R\Hhom_{\Bb} (M,L))} \ar@{<->}[d]\\
\scriptscriptstyle{\Hom(f'^{\ast}M\otimes_{\Aa}g'^{\ast}K,
Rf'^{!}L)} \ar@{<->}[r]_{\scriptscriptstyle{\rm adj}}& \scriptscriptstyle{\Hom(Rf'^{\cdot}_!(f'^{\ast}M\otimes^{L}_\Aa
g'^{\ast}K),L)}
\ar[r]^{\text{\ding{173}}} &
\scriptscriptstyle{\Hom_{\Bb}(M\mathop{\otimes}^{\LL}_\Aa g^{\ast}Rf_!K, L)}.
}
$$

Dans\pageoriginale ce diagramme, \ding{173} (déduit de la flèche
\hyperrefext[XVII][XVII.eq5.2.2.1]{5.2.2.1}) est un isomorphisme (\hyperrefext[XVII][XVII.thm5.2.6]{5.2.6}). La
flèche \ding{172} 
est donc un isomorphisme quel que soit $K$, et ceci implique que
\eqhyperref[XVIII.eq3.1.11.4]{3.1.11.4} soit un isomorphisme.

Le théorème de changement de base pour un morphisme propre se
transpose donc en la 

\begin{enonce*}{Proposition 3.1.12}\etiquette[3.1.12]{XVIII.prop3.1.12}
La flèche \eqhyperref[XVIII.eq3.1.11.4]{3.1.11.4} est un isomorphisme.
\end{enonce*}

On appelle cet isomorphisme \emph{l'isomorphisme d'induction}.

En voici 3 cas particuliers.

\begin{enonce*}{Corollaire 3.1.12.1}\etiquette[3.1.12.1]{XVIII.cor3.1.12.1}
Soient $f:X\rightarrow S$ un morphisme compactifiable, $\mathcal{A}$
et $\mathcal{B}$ deux faisceaux d'anneaux de torsion sur $S$, et
$\varphi:\mathcal{A}\rightarrow \mathcal{B}$ un
homomorphisme. Désignons par $\rho$ la restriction des scalaires de
$\mathcal{B}$ à $\mathcal{A}$. On a alors, pour $K\in \Ob
D^{+}(S,\mathcal{B})$
$$
\rho Rf^{!}K\xrightarrow{\sim} Rf^{!}(\rho K){\quoi}.
$$
\end{enonce*}

Il suffit dans \eqhyperref[XVIII.eq3.1.11.3]{3.1.11.3} de prendre $S=S'$, $X=X'$,
$\mathcal{A}=\mathcal{A}$, $\mathcal{B}=\mathcal{B}$, $M=\mathcal{B}$.

On voit donc que le faisceau d'anneaux joue un rôle bidon dans la
construction de $Rf^{!}$.

\begin{enonce*}{Corollaire 3.1.12.2}\etiquette[3.1.12.2]{XVIII.cor3.1.12.2}
Soient $f:X\rightarrow S$ un morphisme compactifiable, $\mathcal{A}$
un faisceau d'anneaux de torsion sur $S$ et $K\in \Ob D^{-}$
$(S,\mathcal{A})$, $L\in \Ob D^{+}(S,\mathcal{A})$. On a la \emph{formule
d'induction}
\begin{gather*}
R\Hhom(f^{\ast}K,Rf^{!}L)\xrightarrow{\sim} Rf^{!}R\Hhom(K,L)\\
(\text{\rm faire } S=S', X=X'){\quoi}.
\end{gather*}
\end{enonce*}

\begin{enonce*}{Corollaire 3.1.12.3}\etiquette[3.1.12.3]{XVIII.cor3.1.12.3}
Pour tout carré \hyperref[XVIII.subsubsec3.1.11]{3.1.11}, et tout $L\in \Ob D^{+}$
$(S',\mathcal{A})$, on a
\begin{gather*}
Rg'_{\ast}R{f'}^{!}L\xrightarrow{\sim} Rf^{!}Rg_{\ast}L\\
(\text{\rm faire } \mathcal{B}=f^{\ast}\mathcal{A}=M){\quoi}. 
\end{gather*}
\end{enonce*}

\setcounter{subsubsection}{12}
\subsubsection{}\etiquette[3.1.13]{XVIII.subsubsec3.1.13}
Pour\pageoriginale un composé $gh$ de morphismes $S$-compactifiables
$$
X\xrightarrow{h}Y\xrightarrow{g}Z
$$
l'isomorphisme de composition
$$
c_{g,h}:Rg_{!}Rh_{!}\xrightarrow{\sim} R(gh)_{!}
$$
se transpose en un morphisme de composition\nde{Le morphisme transposé va
en réalité dans l'autre sens, mais cela ne pose pas de problème puisque
c'est un isomorphisme.}
\begin{equation*}
c^{!}_{g,h}:Rh^{!}Rg^{!}\xrightarrow{\sim} R(gh)^{!}\tag{3.1.13.1}
\end{equation*}\etiquette[3.1.13.1]{XVIII.eq3.1.13.1}%
vérifiant la condition de cocycle habituelle pour un composé
triple. En d'autres termes, les \emph{catégories $D^{+}$
$(X,f^{\ast}\mathcal{A})$ forment les fibres d'une catégorie fibrée et
cofibrée sur la catégorie d'objets les $S$-schémas} $f:X\rightarrow
S$, ayant pour flèches les morphismes $S$-compactifiables, \emph{les
foncteurs image directe} (\resp \emph{image réciproque}) \emph{étant
les foncteurs} $Rf_{!}$ (\resp $Rf^{!}$). Pour tout carré commutatif
de morphismes compactifiables,
$$
\xymatrix@=1.75cm{
X' \ar[r]_{g} \ar[d]^{f'}& X \ar[d]^{f}\\
S' \ar[r]_{g'}& S,
}
$$
on dispose donc d'un morphisme de cochangement de base (\hyperrefext[XVII][XVII.prop2.1.3]{2.1.3})
\begin{equation*}
ch^{!}:Rf'_{!}Rg^{!}\longrightarrow Rg^{!} Rf_{!}{\quoi}.\tag{3.1.13.2}
\end{equation*}\etiquette[3.1.13.2]{XVIII.eq3.1.13.2}

\subsubsection{}\etiquette[3.1.14]{XVIII.subsubsec3.1.14}
Soit un carré cartésien
$$
\xymatrix@=1.75cm{
X' \ar[r]^{g'} \ar[d]^{f'}& X \ar[d]^{f}\\
S' \ar[r]^{g}& S,
}
$$
avec $f$ compactifiable, et soit $\mathcal{A}$ un faisceau d'anneaux
de torsion sur $S$. On annèle $S'$, $X$ et $X'$ par les images
réciproques de $\mathcal{A}$, encore désignées\pageoriginale 
par $\mathcal{A}$ par abus de notation. L'isomorphisme de changement de base
\begin{equation*}
g^{\ast}Rf_{!}\xrightarrow{\sim} Rf'_{!}{g'}^{\ast}\tag{3.1.14.1}
\end{equation*}\etiquette[3.1.14.1]{XVIII.eq3.1.14.1}%
permet de considérer les catégories $D(S,\mathcal{A})$,
$D(X,\mathcal{A})$, $D(S',\mathcal{A})$ et $D(X',\mathcal{A})$ comme
les fibres d'une catégorie cofibrée sur la catégorie du diagramme
commutatif
$$
\xymatrix@C=2cm@R=1cm{
X \ar[d]& X\ar[l]\ar[d]\\
X & X\ar[l],
}
$$
les foncteurs images directes étant les foncteurs $Rf_{!}$, $Rf'_{!}$,
$g^{\ast}$ et ${g'}^{\ast}$ (sic). Lorsqu'on se restreint aux
catégories $D^{+}$, on obtient même une catégorie fibrée et cofibrée
sur ce diagramme, les foncteurs images directes étant les foncteurs
$Rf^{!}$, $R{f'}^{!}$, $Rg_{\ast}$ et $Rg'_{\ast}$ (sic). On en
déduit
\begin{enumerate}
\renewcommand{\theenumi}{\alph{enumi}}
\renewcommand{\labelenumi}{\theenumi)}
\item un « isomorphisme de transitivité »  $(\sic)$ pour les
foncteurs « images inverses »
$$
Rg_{\ast} Rf^{!}\xrightarrow{\sim} Rf^{!}Rg_{\ast}
$$
déjà obtenu en \hyperref[XVIII.cor3.1.12.3]{3.1.12.3} (avec la même définition par
transposition de \eqhyperref[XVIII.eq3.1.14.1]{3.1.14.1}.

\item un « morphisme de changement de base »  (\hyperrefext[XVII][XVII.prop2.1.3]{2.1.3})
\begin{equation*}
{g'}^{\ast}Rf^{!}\longrightarrow R{f'}^{!}g^{\ast}\tag{3.1.14.2}
\end{equation*}\etiquette[3.1.14.2]{XVIII.eq3.1.14.2}%
qui, d'après \loccit, peut se décrire des deux façons suivantes:
\end{enumerate}

\begin{enumerate}
\renewcommand{\theenumi}{\alph{enumi}}
\renewcommand{\labelenumi}{\theenumi)}
\item par adjonction, se donner un morphisme \eqhyperref[XVIII.eq3.1.14.2]{3.1.14.2} revient à se
donner un morphisme
$$
Rf'_{!}{g'}^{\ast}Rf^{!}\longrightarrow g^{\ast}{\quoi}.
$$
Le théorème de changement de base fournit un tel morphisme comme
composé
$$
Rf'_{!}{g'}^{\ast}Rf^{!}\xleftarrow{\sim}
g^{\ast}Rf_{!}Rf^{!}\xrightarrow{g^{\ast}\ast \text{adj.}} g^{\ast}{\quoi}.
$$

\item par\pageoriginale adjonction, se donner un morphisme \eqhyperref[XVIII.eq3.1.14.2]{3.1.14.2}
revient à se donner un morphisme
$$
Rf^{!}\longrightarrow Rg'_{\ast}R{f'}^{!}g^{\ast}{\quoi}.
$$
L'isomorphisme (\hyperref[XVIII.cor3.1.12.3]{3.1.12.3}) fournit un tel morphisme comme composé
$$
Rf^{!}\xrightarrow[Rf^{!}\ast \text{adj.}]{} Rf^{!}Rg_{\ast}
g^{\ast}\xleftarrow{\sim} Rg'_{\ast}R{f'}^{!}g^{\ast}{\quoi}.
$$
\end{enumerate}

D'après leurs interprétations en terme de catégorie fibrée et
cofibrée, il est clair que \emph{les morphismes} (\hyperref[XVIII.cor3.1.12.3]{3.1.12.3}),
\eqhyperref[XVIII.eq3.1.13.2]{3.1.13.2} \emph{et} \eqhyperref[XVIII.eq3.1.14.2]{3.1.14.2} \emph{vérifient chacun deux
compatibilités du type} \hyperrefext[XII][XII.prop4.4]{4.4} \emph{pour $f$ ou $g$ composé de deux
morphismes\nde{Ces constructions peuvent aussi être interprétées dans le
contexte des structures d'échanges et des foncteurs croisés, \cf \cite{0}.}.}

Les résultats de topologie générale qui suivent seront utilisés au
nᵒ suivant.

\subsubsection{}\etiquette[3.1.15]{XVIII.subsubsec3.1.15}
Soit $f:X\rightarrow S$ un morphisme de sites. Désignons par
$\Gamma(f)$ le site suivant
\begin{enumerate}
\renewcommand{\theenumi}{\alph{enumi}}
\renewcommand{\labelenumi}{\theenumi)}
\item La catégorie $\Gamma(f)$ est la catégorie des triples
$(U,V,\varphi)$ avec
$$
U\in \Ob (X), V\in \Ob S,\quad \varphi\in \Hom
(U,f^{\ast}V){\quoi}.
$$

\item Une famille de morphismes
$(u_{i},v_{i}):(U_{i},V_{i},\varphi_{i})\rightarrow (U,V,\varphi)$ est
couvrante si, dans $X$, les morphismes $u_{i}:U_{i}\rightarrow U$
couvrent $U$.
\end{enumerate}

On a alors

\paragraph{}\etiquette[3.1.15.1]{XVIII.par3.1.15.1}
Le foncteur de $\Gamma(f)$ dans $X$ donné par
$$
(U,V,\varphi)\longmapsto U
$$
est une équivalence de sites $\gamma:X\rightarrow \Gamma(f)$.

\paragraph{}\etiquette[3.1.15.2]{XVIII.par3.1.15.2}
Le foncteur $f^{\ast}_{\Gamma}$ de $S$ dans $\Gamma(f)$
$$
V\longmapsto (f^{\ast}V,V,\Id_{f^{\ast}V})
$$
est un morphisme de site, noté $f_{\Gamma}:\Gamma(f)\rightarrow S$, ou
simplement $f$ par abus de notation. Il vérifie
$f_{\Gamma}\gamma\simeq f$. Le foncteur $f^{\ast}_{\Gamma}$ admet un
adjoint à gauche $f_{\Gamma\bigdot}$:
$$
f_{\Gamma\bigdot}:(U,V,\varphi)\longmapsto V{\quoi}.
$$
Le\pageoriginale foncteur $f_{\Gamma\bigdot}$ est donc un comorphisme
de site de $\Gamma(f)$ dans $S$\nde{Cette terminologie n'étant pas
familière, il convient de préciser ici qu'il s'agit des notions
intervenant dans \hyperrefext[III][III.sec2]{2} et \hyperrefext[IV][IV.subsec4.7]{4.7} (on prendra garde au
fait que dans l'édition originale, à l'endroit de cette deuxième référence,
des coquilles empêchaient de comprendre
dans quel sens allaient au juste les flèches).},
et définit le même morphisme de topos
que le morphisme de site $f_{\Gamma}$.
\begin{equation*}
\vcenter{
\xymatrix@R=2.2cm@C=1.75cm{
X &         &\Gamma(f) \ar[ll]^{\approx}
\ar@<4pt>[ld]^{f_{\Gamma}.}& \\
  &    S \ar[lu]^{f^{\ast}}  \ar@<4pt>[ru]^{f^{\ast}_{\Gamma}} & .
}}\tag{3.1.15.3}
\end{equation*}\etiquette[3.1.15.3]{XVIII.eq3.1.15.3}

On en déduit la règle suivante pour calculer l'image inverse d'un
faisceau\nde{Il est ici fait usage de \hyperrefext[III][III.prop2.3]{2.3} ou de
\hyperrefext[IV][IV.eq4.7.3]{4.7.3}.}

\setcounter{paragraph}{3}
\paragraph{}\etiquette[3.1.15.4]{XVIII.par3.1.15.4}
Soient $\Ff$ un faisceau sur $S$,
$f^{\bigdot}_{\Gamma}\Ff$ le préfaisceau sur $\Gamma(f)$ donné
par
$$
f^{\bigdot}_{\Gamma}\Ff(U,V,\varphi)=\Ff(V){\quoi},
$$
et $af^{\bigdot}_{\Gamma}\Ff$ le faisceau engendré. Alors
$$
f^{\ast}\Ff(U)\simeq
af^{\bigdot}_{\Gamma}\Ff(U,V,\varphi){\quoi}.
$$

\subsubsection{}\etiquette[3.1.16]{XVIII.subsubsec3.1.16}
Soient $f:X\rightarrow S$ un morphisme de sites, $\mathcal{A}$ un
faisceau d'anneaux sur $S$, $\mathcal{B}$ un faisceau d'anneaux sur
$X$ et $d$ un entier. Soit un objet $K$ du type suivant

\paragraph{}\etiquette[3.1.16.1]{XVIII.par3.1.16.1}
$K$ associe à chaque objet $U$ de $X$ un complexe $K(U)$ de
$(\mathcal{A},f_{|U\ast}(\mathcal{B}|U))$-modules sur $S$, nul en degré
$>d$, et $K(U)$ est un foncteur \emph{covariant} de $U$.

Pour tout faisceau de $\mathcal{A}$-modules $ℱ$ sur $S$, on
désignera par $K^{!}ℱ$ le complexe de faisceaux sur $X$ engendré
par le complexe de préfaisceaux de $\mathcal{B}$-modules:
$$
U\longmapsto \Hom_{\mathcal{A}}(K(U),ℱ){\quoi}.
$$

On désignera par $RK^{!}$ le dérivé de $K^{!}$
$$
RK^{!}:D^{+}(S,\mathcal{A})\longrightarrow
D^{+}(X,\mathcal{B}){\quoi}.
$$
Si $U:K_{1}\rightarrow K_{2}$ est un morphisme d'objets \eqhyperref[XVIII.eq3.1.11.1]{3.1.11.1},
alors $U$ induit\pageoriginale des morphismes de foncteurs
$$
{}^{t}u:K^{!}_{2}\longrightarrow K^!_{1}\quaet RK^{!}_{2}\longrightarrow
RK^{!}_{1}{\quoi}.
$$

Sous les hypothèses précédentes, on a

\begin{enonce*}{Proposition 3.1.17}\etiquette[3.1.17]{XVIII.prop3.1.17}
\begin{enumerate}
\renewcommand{\theenumi}{\roman{enumi}}
\renewcommand{\labelenumi}{(\theenumi)}
\item Pour tout point $x$ de $X$ et pour tout $L\in
\Ob D^{+}(S,\mathcal{A})$, on a, désignant par $V(x)$ la catégorie des
voisinages de $x$
\begin{equation*}
ℋ^{n}(RK^{!}L)_{x}\simeq \varinjlim_{U\in \Ob V(x)}
\Ext^{n}(K(U),L){\quoi}.\tag{3.1.17.1}
\end{equation*}\etiquette[3.1.17.1]{XVIII.eq3.1.17.1}

\item On a donc une suite spectrale
\begin{equation*}
E^{pq}_{2}=\varinjlim_{U\in \Ob V(x)}
\Ext^{p}(ℋ^{-q}(K(U)),L)\Longrightarrow
ℋ^{n}(K^{!}L)_{x}{\quoi}.\tag{3.1.17.2}
\end{equation*}\etiquette[3.1.17.2]{XVIII.eq3.1.17.2}

\item Si $X^{0}$ est un ensemble conservatif de points de $X$, et
$u:K_{1}\rightarrow K_{2}$ un morphisme d'objets \eqhyperref[XVIII.eq3.1.11.1]{3.1.11.1}, alors,
pour que
$$
{}^{t}u:RK^{!}_{2}\longrightarrow RK^{!}_{1}
$$
soit un isomorphisme, il faut et il suffit que pour tout $n$ et tout
$x\in X^{0}$, le morphisme de pro-objets de $\Mod(S,\mathcal{A})$
\begin{equation*}
u_{x}:\mathop{``\varprojlim\mbox{''}}_{U\in \Ob V(x)}
ℋ^{n}K_{1}(U)\longrightarrow \mathop{``\varprojlim\mbox{''}}_{U\in
\Ob V(x)} ℋ^{n}K_{2}(U)\tag{3.1.17.3}
\end{equation*}\etiquette[3.1.17.3]{XVIII.eq3.1.17.3}%
soit un isomorphisme.
\end{enumerate}
\end{enonce*}

Pour prouver (i), on prend pour $L$ un complexe d'injectifs, de sorte
que \eqhyperref[XVIII.eq3.1.17.1]{3.1.17.1} est trivial; il est clair que (i) $\Rightarrow$ (ii).

Que la condition \eqhyperref[XVIII.eq3.1.17.3]{3.1.17.3} de (iii) soit suffisante résulte aussitôt
de \eqhyperref[XVIII.eq3.1.17.2]{3.1.17.2}. Nous n'aurons pas à faire usage de ce qu'elle est
nécessaire. Si ${}^{t}u$ est un isomorphisme, alors, d'après
(\hyperref[XVIII.par3.1.16.1]{3.1.16.1}), pour tout faisceau injectif $I$
sur $(S,\mathcal{A})$, on a
$$
\Hom(\mathop{``\varprojlim\mbox{''}}_{V(x)}ℋ^{n}K_{2}(U),I)\xrightarrow{\sim}\Hom(\mathop{``\varprojlim\mbox{''}}_{V(x)}
ℋ^{n}K_{1}(U),I){\quoi}, 
$$
et on conclut en notant que dans toute catégorie abélienne ayant assez
d'injectifs (ici, $\Mod(S,\mathcal{A})$), les foncteurs $\Hom(\ast,I)$
pour $I$ injectif de $\mathcal{A}$ forment un système conservatif de
foncteurs sur $\Pro(\mathcal{A})$.

\begin{enonce*}[remark]{Exemple 3.1.18}\etiquette[3.1.18]{XVIII.exem3.1.18}
Soient\pageoriginale $f:X\rightarrow S$ un morphisme compactifiable de
schémas, et $\mathcal{A}$ un faisceau d'anneaux de torsion sur $S$. On
annèle $X$ par $f^{\ast}\mathcal{A}$. Choisissons un entier $d$ qui
majore la dimension des fibres de $f$, ainsi qu'une compactification
de $f$, et posons \eqhyperref[XVIII.eq3.1.4.7]{3.1.4.7})
\begin{equation*}
K(U)=f^{\bigdot}_{!}((f^{\ast}\mathcal{A})_{U}){\quoi}.\tag{3.1.18.1}
\end{equation*}\etiquette[3.1.18.1]{XVIII.eq3.1.18.1}
\end{enonce*}

D'après la formule \eqhyperref[XVIII.eq3.1.3.1]{3.1.3.1} pour l'adjoint de $f^{\bigdot}_{!}$, on a
un isomorphisme
\begin{equation*}
f^{!}_{\bigdot}\simeq K^{!}:\Mod(S,\mathcal{A})\longrightarrow C^{b}\Mod(X,f^{\ast}\mathcal{A}){\quoi}.\tag{3.1.18.2}
\end{equation*}\etiquette[3.1.18.2]{XVIII.eq3.1.18.2}%
et donc
\begin{equation*}
Rf^{!}=RK^{!}:D^{+}(S,\mathcal{A})\longrightarrow
D^{+}(X,f^{\ast}\mathcal{A}){\quoi}.\tag{3.1.18.3}
\end{equation*}\etiquette[3.1.18.3]{XVIII.eq3.1.18.3}

\begin{enonce*}[remark]{Exemple 3.1.19}\etiquette[3.1.19]{XVIII.exem3.1.19}
Soient $f:X\rightarrow S$, $\mathcal{A}$, $\mathcal{B}$ et $d$ comme
en \hyperref[XVIII.subsubsec3.1.16]{3.1.16}. Soit $K$ un objet \sisi{de}{du}
type suivant.
\end{enonce*}

\setcounter{subsubsection}{19}
\setcounter{paragraph}{0}
\paragraph{}\etiquette[3.1.19.1]{XVIII.par3.1.19.1}
Pour tout $W=(U,V,\varphi)\in \Ob \Gamma(f)$
(\hyperref[XVIII.subsubsec3.1.15]{3.1.15})\nde{Par définition, $\varphi$ est un
$X$-morphisme $U\to V\times_S X$. Pour bien comprendre cet exemple, il faut
prendre le point de vue équivalent consistant à interpréter $\varphi$ comme
un $S$-morphisme $\varphi\colon U\to V$.}, $K(W)$ est un 
complexe de $(\mathcal{A}_{V},\varphi_{\ast}B_{|U})$-Modules sur $V$,
fonctoriel en $W$: pour tout morphisme $(u,v):
W_{1}=(U_{1},V_{1},\varphi_{1})\to W_{2}=(U_{2},V_{2},\varphi_{2})$ il
est donné
\begin{align*}
K(u,v) &: K(W_{1})\longrightarrow v^{\ast}K(W_{2}),\quad\mbox{\ie}\\
K(u,v) &: v_{!}K(W_{1})\longrightarrow K(W_{2}){\quoi}.
\end{align*}

À tout objet $K$ (\hyperref[XVIII.par3.1.19.1]{3.1.19.1}) est associé un objet $K'$
(\hyperref[XVIII.par3.1.16.1]{3.1.16.1}) 
relatif à $f_{\Gamma}:\Gamma(f)\rightarrow S$, donné par
\begin{equation*}
K'(U,V,\varphi)=j_{!}K(U,V,\varphi)\tag{3.1.19.2}
\end{equation*}\etiquette[3.1.19.2]{XVIII.eq3.1.19.2}%
pour $j$ « morphisme d'inclusion »  de $V$ dans $S$. On pose
$$
K^{!}\mathop{=}_{\defn} {K'}^{!}\quaet RK^{!}\mathop{=}_{\defn}
R{K'}^{!}{\quoi}.
$$

\begin{enonce*}[remark]{Exemple 3.1.20}\etiquette[3.1.20]{XVIII.exem3.1.20}
Avec les notations de \hyperref[XVIII.exem3.1.19]{3.1.19}, faisons
$\mathcal{B}=f^{\ast}\mathcal{A}$ et posons
\begin{equation*}
K(U,V,\varphi)=\mathcal{A}_{V}\quad\text{\rm(en degré 0)}\tag{3.1.20.1}
\end{equation*}\etiquette[3.1.20.1]{XVIII.eq3.1.20.1}
\end{enonce*}

D'après\pageoriginale \hyperref[XVIII.par3.1.15.4]{3.1.15.4}, on a alors des
isomorphismes de foncteurs
\begin{gather*}
K^{!}\simeq f^{\ast}\tag{3.1.20.2}\\
RK^{!}\simeq f^{\ast}{\quoi}.\tag{3.1.20.3}
\end{gather*}\etiquette[3.1.20.2]{XVIII.eq3.1.20.2}
\etiquette[3.1.20.3]{XVIII.eq3.1.20.3}

\subsection{Dualité de Poincaré}\etiquette[3.2]{XVIII.subsec3.2}

\subsubsection{}\etiquette[3.2.1]{XVIII.subsubsec3.2.1}
Soit $f$ un morphisme plat compactifiable purement de dimension
relative $d$ (ou, plus généralement, vérifiant la condition
$(\ast)_{d}$ de \hyperref[XVIII.thm2.9]{2.9}):
$$
\xymatrix@C=2cm@R=1.5cm@M=.2cm{
X \ar@{^{(}->}[r]^{j} \ar[d]_{f}& \bar{X}\ar[ld]_{\bar{f}} &\\
S &.
}
$$
Soit $n$ un entier $\geq 1$ inversible sur $S$ et supposons $S$ annelé
par un faisceau d'anneaux $\mathcal{A}$ tel que
$n\mathcal{A}=0$. D'après \hyperref[XVIII.exem3.1.18]{3.1.18}, le foncteur
$Rf^{!}:D^{+}(S,\mathcal{A})\rightarrow D^{+}(X,\mathcal{A})$ peut se
calculer par le procédé de \hyperref[XVIII.subsubsec3.1.16]{3.1.16}, pour $K$ donné
par \eqhyperref[XVIII.eq3.1.18.1]{3.1.18.1}. 

Soit $K'$ le foncteur du type (\hyperref[XVIII.par3.1.19.1]{3.1.19.1}) qui à
$(U,V,\varphi)\in \Ob \Gamma(f)$ associe le complexe de faisceaux
$\ZZ/n(-d)[-2d]$ 
sur $V$, réduit à $\ZZ/n(-d)$ placé en degré $2d$, et soit $K''$
l'objet du type (\hyperref[XVIII.par3.1.16.1]{3.1.16.1}) correspondant
\eqhyperref[XVIII.eq3.1.19.2]{3.1.19.2}, relatif à $f:\Gamma(f)\rightarrow S$.

On peut regarder $K$ et $K''$ comme des objets
\hyperref[XVIII.par3.1.16.1]{3.1.16.1} relatif à 
$\Gamma(f)\rightarrow S$. De plus, pour tout $(U,V,\varphi)\in \Ob
\Gamma(f)$, le morphisme trace \hyperref[XVIII.thm2.9]{2.9} définit un morphisme
de faisceaux sur $V$
$$
\Tr:R^{2d}\varphi_{!}\ZZ/n\longrightarrow
\ZZ/n(-d){\quoi},
$$
d'où, désignant par $j$ le morphisme de $V$ dans $S$, un morphisme de
foncteurs
\begin{gather*}
K(U)=\bar{f}_{\ast}\tau_{\leq
2d}C^{\ast}_{\ell}\ZZ/n_{U}\longrightarrow
ℋ^{2d}(K(U))[-2d]\simeq\tag{3.2.1.1}\\
R^{2d}f_{!}(\ZZ/n)_{U}\simeq j_{!}R^{2d}\varphi ! \ZZ/n
\xrightarrow{\Tr_{\varphi}} j_{!}
\ZZ/N(-d)=K''(U,V,\varphi)){\quoi}.
\end{gather*}\etiquette[3.2.1.1]{XVIII.eq3.2.1.1}
 
Une\pageoriginale forme tordue de \hyperref[XVIII.exem3.1.20]{3.1.20} montre que
$$
{RK''}^{!}=Rf^{\ast}(d)[2d]{\quoi}.
$$
D'après \hyperref[XVIII.subsubsec3.1.16]{3.1.16}, on dispose donc d'un morphisme composé
\begin{equation*}
t_{f}:Rf^{\ast}(d)[2d]\simeq {RK''}^{!}\longrightarrow Rk^{!}\simeq
Rf^{!}{\quoi}.\tag{3.2.1.2}
\end{equation*}\etiquette[3.2.1.2]{XVIII.eq3.2.1.2}

\begin{enonce*}[remark]{Remarque 3.2.2}\etiquette[3.2.2]{XVIII.rem3.2.2}
Si le foncteur $Rf^{!}$ est de dimension cohomologique finie, alors
$t_f$ est encore défini en tant que morphisme entre foncteurs de
$D(S,\mathcal{A})$ dans $D(X,f^{\ast}\mathcal{A})$, (définition en
\hyperref[XVIII.rem3.1.5]{3.1.5}).
\end{enonce*}

On dispose d'une autre méthode pour construire une flèche \eqhyperref[XVIII.eq3.2.1.2]{3.2.1.2},
par adjonction à partir de \eqhyperref[XVIII.eq2.13.2]{2.13.2}
$$
\Tr:Rf_{!}(f^{\ast}L(d)[2d])\longrightarrow L{\quoi}.
$$
Ces deux flèches coïncident:

\begin{enonce*}{Lemme 3.2.3}\etiquette[3.2.3]{XVIII.lem3.2.3}
Pour $L\in \Ob D^{+}(S,\mathcal{A})$ (\resp $D(S,\mathcal{A})$ \emph{si
le foncteur $Rf^{!}$ est de dimension cohomologique finie}), le
diagramme
$$
\xymatrix@C=2.5cm@R=1.75cm{
Rf_!(f^{\ast}L(d)[2d]) \ar[d]^{\Tr_f} \ar[r]^{Rf_!(t_f)}  & Rf_!
Rf^{!}L \ar[d]^{\txt{\rm adjonction}}\\
L                              \ar@{=}[r] & L
}
$$
est commutatif. 
\end{enonce*}

Donnons tout d'abord une variante de la définition de $t_f$, en termes
du foncteur $f^!_\bigdot$ \eqhyperref[XVIII.eq3.1.4.10]{3.1.4.10}.
Soient $U\in \Ob X_{\text{\rm et}}$,
$V\in \Ob S_{\text{\rm et}}$ et $\varphi\in \Hom_{f}(U,V)$. Par
définition,
$$
f^{!}_{\bigdot}(L)(U)=\Hom(f^{\bigdot}_{!}\mathcal{A}_{U},L){\quoi}\;;
$$
le morphisme trace définit un morphisme
$$
\Tr_{f}:f^{\bigdot}_{!}\mathcal{A}_{U}\longrightarrow
\mathcal{A}_{V}(-d)[-2d]
$$
qui se transpose en
$$
t_{\varphi}:f^{!}_{\bigdot}(L)(U)\longleftarrow L(d)[2d](V){\quoi}.
$$
Les\pageoriginale morphismes $t_{\varphi}$ définissent un morphisme
$$
t_{f}:f^{\ast}L(d)[2d]\longrightarrow f^{!}_{\bigdot}L
$$
dont on laisse au lecteur le soin de vérifier que, dans la catégorie
dérivée, il coïncide avec \eqhyperref[XVIII.eq3.2.1.2]{3.2.1.2}. Cette description montre que les
diagrammes suivants sont commutatifs au niveau des complexes, pour $L$
complexe de Modules:
$$
\xymatrix@R=1.25cm{
\Hom^{\bigdot}(\Aa_V,L) \ar[r] \ar[d]^{\Tr_f}& \Hom^{\bigdot}(f^{\ast}\Aa_V(d) [2d], f^{\ast}
L(d) [2d])\ar[d]^{t_f}\\
\Hom(f^{\bigdot}_!f^{\ast}\Aa_V(d)[2d],L) \ar@{-}[r]^{\sim}& \Hom^{\bigdot}(f^{\ast}\Aa_V(d)[2d], f^{!}_{\bigdot}L).
}
$$
Par fonctorialité des flèches de ce diagramme en $\mathcal{A}_{V}$, on
déduit que pour tout faisceau, ou tout complexe de faisceaux $K$, le
diagramme suivant est commutatif:
$$
\xymatrix{
\Hom^{\bigdot}(K,L) \ar[r] \ar[d]^{\Tr_f}& \Hom^{\bigdot}(f^{\ast}K(d)[2d],
f^{\ast}L(d)[2d]) \ar[d]\\
\Hom^{\bigdot}(f^{\bigdot}_!f^{\ast}K(d)[2d],L) \ar@{<->}[r]^{\sim} &\Hom^{\bigdot}(f^{\ast}K(d)[2d], f^{!}_{\bigdot}L).
}
$$
Prenant $K=L$ et $\Id:L\rightarrow L$, on trouve \hyperref[XVIII.lem3.2.3]{3.2.3}
au niveau des complexes.

Je serais reconnaissant à toute personne ayant compris cette
démonstration de me l'expliquer.

\setcounter{subsubsection}{3}
\subsubsection{}\etiquette[3.2.4]{XVIII.subsubsec3.2.4}
On déduit aussitôt de ce lemme et de la définition
(\hyperref[XVIII.subsubsec3.1.13]{3.1.13}) des 
isomorphismes de composition par adjonction que si $f$ est le composé
de deux morphismes plats de présentation finie compactifiables et
purement de dimension relative $d_{1}$ et $d_{2}:f=gh$, et si
$d=d_{1}+d_{2}$, alors le diagramme
$$
\xymatrix{
h^{\ast}(g^{\ast}L(d_2)[2d_2])(d_1)[2d_1])
\ar@{-}[r]^(.6){\sim}
\ar[d]^{t_h\ast t_g}&(gh)^{\ast}L(d)[2d] \ar[d]^{t_{gh}}\\
Rh^{!} \; Rg^{!} \; L  \ar@{-}[r]^{\sim}&  R(gh)^{!}L
}
$$
est\pageoriginale commutatif.

\begin{enonce*}{Théorème 3.2.5}[Dualité de Poincaré]%
\etiquette[3.2.5]{XVIII.thm3.2.5}
Soient $f:X\rightarrow S$ un morphisme
lisse compactifiable \emph{($S$ est donc quasi-compact
quasi-séparé (sic))}, $d$ la fonction localement constante sur $X$
« dimension relative de $f$ », $n\geq 1$ un entier inversible sur
$S$, et $\mathcal{A}$ un faisceau d'anneaux sur $S$ vérifiant
$n\mathcal{A}=0$. Alors le foncteur
$$
Rf_{!}:D(X,f^{\ast}\mathcal{A})\longrightarrow D(S,\mathcal{A})
$$
admet pour adjoint à droite le foncteur de $D(S,\mathcal{A})$ dans
$D(X,f^{\ast}\mathcal{A})$:
$$
K\longmapsto f^{\ast}K(d)[2d]
$$
avec le morphisme \eqhyperref[XVIII.eq2.13.2]{2.13.2}
$$
\Tr_{f}:Rf_{!}f^{\ast}K(d)[2d]\longrightarrow K
$$
pour flèche d'adjonction:
$$
\Hom(K,f^{\ast}L(d)[2d])\xrightarrow{\sim} \Hom(Rf_{!}K,L){\quoi}.
$$
\end{enonce*}

\noindent
\emph{Démonstration}. On se ramène au cas $d$ constant. D'après la
description \hyperref[XVIII.subsubsec3.2.1]{3.2.1} de la flèche $t_{f}$
\eqhyperref[XVIII.eq3.2.1.2]{3.2.1.2}, le théorème \hyperref[XVIII.thm2.14]{2.14}, 
joint au critère \hyperref[XVIII.prop3.1.17]{3.1.17} (iii), implique que la flèche $t_{f}$ est un
isomorphisme (pour $L\in D^{+}(S,\mathcal{A})$)
$$
t_{f}:f^{\ast}L(d)[2d]\xrightarrow{\sim} Rf^{!}L{\quoi}.
$$

En particulier, le foncteur $Rf^{!}$ est de dimension cohomologique
finie, donc est défini sur la catégorie dérivée entière. Le morphisme
$t_{f}$ reste un isomorphisme pour $L\in D(S,\mathcal{A})$\nde{Plus
précisément, le fait
que $Rf^!$ soit de dimension cohomologique finie permet de montrer que l'on
peut définir le foncteur $Rf^!$ sur la catégorie dérivée entière
de façon à ce
que pour un complexe $I^\bigdot$ de faisceaux injectifs, $Rf^!I^\bigdot$
soit le complexe simple associé au complexe double $f^!_\bigdot I^\bigdot$.
Dès lors, pour tout $n\in\ZZ$, le morphisme induit au niveau du
$n$-ième faisceau de cohomologie par $t_f$ évalué sur $I^\bigdot$
s'identifie au même morphisme pour un tronqué bête $I^{\geq k}$ pour
$k$ assez petit, ce qui permet de conclure.}.

Par\pageoriginale définition, $Rf^{!}$ est l'adjoint à droite du
foncteur $Rf_{!}$; d'après \hyperref[XVIII.lem3.2.3]{3.2.3}, la flèche
d'adjonction de $Rf_{!}Rf^{!}K$ dans $K$ s'identifie,
via $t_{f}$, à $\Tr_{f}$, et ceci
achève la démonstration de \hyperref[XVIII.thm3.2.5]{3.2.5}.

On identifiera dorénavant à l'aide de $t_{f}$ les foncteurs
$f^{\ast}L(d)[2d]$ et $Rf^{!}L$.

\setcounter{subsubsection}{5}
\subsubsection{}\etiquette[3.2.6]{XVIII.subsubsec3.2.6}
Soient $S$ le spectre d'un corps algébriquement clos, $k$, $X$ un
schéma compactifiable et lisse sur $k$ purement de dimension relative
$d$, $n$ un entier inversible dans $k$ et $F$ un faisceau de
$\ZZ/n$-modules localement constant constructible, \ie un « système de coefficients »  tué par $n$. Désignant par
${}^\vee$ un $\ZZ/n$-dual, on a
$$
R\Hhom(F,\ZZ/n)=F^{\vee}\quad\text{\rm (en degré 0)}{\quoi}.
$$
Cette formule résulte de ce que $\ZZ/n$ est un
$\ZZ/n$-module injectif. L'isomorphisme 3.2.5.1, ou, plutôt
l'isomorphisme qui s'en déduit
$$
R\Gamma\;R\Hhom(F,\ZZ/n(d)[2d])\xrightarrow{\sim}
R\Hom(R\Gamma_{c}(F),\ZZ/n){\quoi},
$$
s'écrit donc ici
\begin{equation*}
R\Gamma (F^{\vee}(d)[2d])\xrightarrow{\sim} R\Hom
(R\Gamma_{c}(F),\ZZ/n){\quoi}.\tag{3.2.6.1}
\end{equation*}\etiquette[3.2.6.1]{XVIII.eq3.2.6.1}%
Passant aux groupes de cohomologie, et utilisant à nouveau que
$\ZZ/n$ est un $\ZZ/n$-module injectif, ceci donne
\begin{equation*}
H^{2d-i}(X,F^{\vee}(d))\xrightarrow{\sim}H^{i}_{c}(X,F)^{\vee}{\quoi},\tag{3.2.6.2}
\end{equation*}\etiquette[3.2.6.2]{XVIII.eq3.2.6.2}%
qui est la forme habituelle de la dualité de Poincaré.

\begin{thebibliography}{99}\pageoriginale%

\sisi{}{%
\bibitem{0} J. \textsc{Ayoub}, Les six opérations de Grothendieck et le
formalisme des cycles évanescents dans le monde motivique. I,
Astérisque 314 (2007).

\bibitem{0bis} B. \textsc{Conrad}, A modern proof of Chevalley's theorem on
algebraic groups, Journal of the Ramanujan Mathematical Society 17 (2002),
p.~1--18.

\bibitem{0ter} D. \textsc{Ferrand}, On the non additivity of the trace
in derived categories, \url{http://arxiv.org/abs/math/0506589}

}

\bibitem{1} J. \textsc{Giraud}, Cohomologie non abélienne. Columbia
  university. 

\bibitem{2} A. \textsc{Grothendieck}, Le groupe de Brauer
  III. dans: dix exposés sur 
la cohomologie des schémas. North Holl. Publ. Co.

\sisi{}%
{\bibitem{2bis} A. \textsc{Grothendieck}, Technique de descente et théorème
d'existence en géométrie algébrique. VI. Les schémas de Picard : propriétés
générales, Séminaire Bourbaki 1961-1962, exposé nᵒ 236,
pages 221-243.}

\sisi{\bibitem{3} L. \textsc{Illusie}, Thèse}%
{\bibitem{3} L. \textsc{Illusie}, Complexe cotangent et déformations II.
Lecture notes in math. 283. Springer Verlag.}

\sisi{}{%
\bibitem{3bis} F. \textsc{Knudsen}, D. \textsc{Mumford}, The
projectivity of the moduli space of stable curves. I : Preliminaries on
``det'' and ``Div'', Mathematica Scandinavica 39 (1976), p.~19--55.

\bibitem{3ter} U. \textsc{Köpf},
Über eigentliche Familien algebraischer Varietäten
über affinoiden Räumen, Schriftenreihe des Mathematischen Instituts
der Universität Münster. 2. Serie, Heft 7 (1974).}

\bibitem{4} D. \textsc{Mumford}, Geometric invariant
  theory. Springer-Verlag 1965. 

\sisi{}{%
\bibitem{4bis} F. \textsc{Orgogozo}, Isomotifs de dimension inférieure ou
égale à un, Manuscripta Mathematica 115 (2004), p.~339--360.}

\bibitem{5} M. \textsc{Raynaud}, Faisceaux amples sur les schémas
  en groupes et les espaces homogènes. Lecture notes in math. 119. Springer Verlag.

\sisi{}{%
\bibitem{5bis} M. \textsc{Raynaud}, Géométrie analytique rigide d'après
Tate, Kiehl. Mémoires de la Société Mathématique de France, tome 39--40
(1974), p.~319--327.}

\bibitem{6} J.P. \textsc{Serre}, Groupes algébriques et corps de
classe. Publ. Inst. Math. Nancago. Hermann 1959.

\bibitem{7} J.L. \textsc{Verdier}, Dualité dans la cohomologie des espaces
localement compacts. Sém. Bourbaki 300. Nov. 1965.

\bibitem{8} J.L. \textsc{Verdier}, A duality theorem in the etale
  cohomology of schemes. Woods-hole seminar on alg. geometry. July 1964.

\bibitem{9} Catégories dérivées, par J.L. \textsc{Verdier}. Notes
miméographiées par l'IHÉS.

\end{thebibliography}


%2.17  - 106
%IX 3.1 - 138
%XVII  6.3.18.2 - 345
%1.2.0 - 855
%1.2.7.3 - 865
%1.2.2.11 - 892
%1.2.12 - 920
%6.3.3.1 - 1026
%1.3.12.4 - 1631, 1636, 1704
%1.3.2.7 - 1669
%1.2.9.7 - 1778
%1.2.16.2 - 1784
%2.12.3 - 4257
%3.1.9.10 -  5337
%3.2.5.1 - 5461

\ifx\danslelivre\undefined
\end{document}
\fi
